\begin{derivation}[单模发射与吸收的有限时间谱核]
在长度为 $T$ 的矩形作用窗中，分别计算发射和吸收一个模式量子的概率，区分玻色占据因子与有限时间能量分辨率的来源。

对单一模式 $\lambda$，相互作用绘景中与一量子交换有关的部分写成
\[
\hat H_I(t)
=\hbar\mathcal G_{fi}^{(\lambda)}
|f\rangle\langle i|\hat a_\lambda^\dagger
 \ee^{\ii\Delta_{\rm em}t}
+\hbar\mathcal G_{fi}^{(\lambda)}
|f\rangle\langle i|\hat a_\lambda
 \ee^{\ii\Delta_{\rm abs}t}
+\mathrm{H.c.},
\]
其中
\[
\Delta_{\rm em}
=\frac{E_f-E_i+\hbar\omega_\lambda}{\hbar},
\qquad
\Delta_{\rm abs}
=\frac{E_f-E_i-\hbar\omega_\lambda}{\hbar}.
\]
一阶 Dyson 振幅对发射过程给
\begin{align*}
\mathcal A_{\rm em}^{(1)}
&=-\frac{\ii}{\hbar}
\int_{-T/2}^{T/2}\dd t\,
\langle f,n_\lambda+1|\hat H_I(t)|i,n_\lambda\rangle\\
&=-\ii\mathcal G_{fi}^{(\lambda)}\sqrt{n_\lambda+1}
\,T\,\sinc\!\left(\frac{\Delta_{\rm em}T}{2}\right).
\end{align*}
这里 $\sqrt{n_\lambda+1}$ 完全来自
$\hat a_\lambda^\dagger|n_\lambda\rangle
=\sqrt{n_\lambda+1}|n_\lambda+1\rangle$。因此
\[
P_{\rm em}^{(1)}(T)
=|\mathcal G_{fi}^{(\lambda)}|^2(n_\lambda+1)T^2
\sinc^2\!\left(\frac{\Delta_{\rm em}T}{2}\right).
\]
同理，吸收过程使用
$\hat a_\lambda|n_\lambda\rangle=\sqrt{n_\lambda}|n_\lambda-1\rangle$，得到
\[
P_{\rm abs}^{(1)}(T)
=|\mathcal G_{fi}^{(\lambda)}|^2n_\lambda T^2
\sinc^2\!\left(\frac{\Delta_{\rm abs}T}{2}\right).
\]
这就是正文\eqnrefs{eq:mode-emission-absorption-finite-time-dp68}。

有限时间核本身已经给出能量分辨率。令
\[
K_T(\Delta)
\equiv\frac{T}{2\pi}\sinc^2\!\left(\frac{\Delta T}{2}\right).
\]
变量代换 $u=\Delta T/2$ 给
\[
\int_{-\infty}^{\infty}\dd\Delta\,K_T(\Delta)
=\frac1\pi\int_{-\infty}^{\infty}\dd u\,\sinc^2u=1.
\]
因此 $K_T$ 是单位面积、宽度随 $T^{-1}$ 缩小的近似恒等核；只有在末态连续且其余因子在该宽度内变化缓慢时，$K_T$ 才在分布极限收缩为 $\delta(\Delta)$。所以 $n_\lambda+1$ 的真空项来自玻色代数，而能量守恒 $\delta$ 来自长时间极限，两者是完全不同的数学来源。
\end{derivation}
